% id: 0182
% book: RPM 중학 3-1 · 02 무리수와 실수
% section: 유형 04 무리수를 수직선 위에 나타내기
% figure: crop:fig-0182.png
% answer: $-1+\sqrt{8}$, $-1-\sqrt{8}$
% answer_source: 답지
% variant_level: 0 (원본 전사)
% 이 파일은 build.mjs 가 items.json 에서 만든다. 고칠 때는 items.json 을 고친다.
\smash{\raisebox{1.5mm}{\dmkichul{RPM 중학 3-1 0182번 · 29쪽 · 대표문제}}}
\noindent\begin{minipage}[t]{0.58\linewidth}\vspace{0pt}\raggedright 오른쪽 그림은 한 눈금의 길이가 1인 모눈종이 위에 수직선과 직각삼각형~$\pt{ABC}$를 그리고, 점~$\pt{A}$를 중심으로 하고 $\seg{AC}$를 반지름으로 하는 원을 그린 것이다. 원과 수직선이 만나는 두 점을 각각 $\pt{P}$, $\pt{Q}$라 할 때, 두 점~$\pt{P}$, $\pt{Q}$에 대응하는 수를 차례대로 구하시오.\end{minipage}\hfill\begin{minipage}[t]{0.4\linewidth}\vspace{0pt}\centering \includegraphics[width=110.5pt]{figures/fig-0182.png}\end{minipage}\par
